% Generated draft: requires mathematical and source review.
\documentclass[11pt]{article}
\usepackage[margin=25mm]{geometry}
\usepackage{fontspec}
\setmainfont{Latin Modern Roman}
\usepackage{amsmath,amssymb,mathtools,unicode-math}
\setmathfont{Latin Modern Math}
\usepackage{enumitem,needspace,xurl,hyperref}
\setlength{\parindent}{0pt}
\setlength{\parskip}{5pt}
\setlength{\emergencystretch}{3em}
\newcommand{\N}{\mathbb{N}}
\newcommand{\Z}{\mathbb{Z}}
\newcommand{\Q}{\mathbb{Q}}
\newcommand{\R}{\mathbb{R}}
\newcommand{\C}{\mathbb{C}}
\newcommand{\F}{\mathbb{F}}
\newcommand{\A}{\mathbb{A}}
\newcommand{\PP}{\mathbb{P}}
\newcommand{\E}{\mathbb{E}}
\newcommand{\eps}{\varepsilon}
\newcommand{\dd}{\,\mathrm{d}}
\begin{document}
\section*{UnsolvedMath ID: 800699}
\textbf{Source problem code:} catalogue entry
\subsection{Definitions and mathematical statement}
Let \(F\) be a field, and let \(\operatorname{Br}(F)\) denote its Brauer group: the group of Brauer equivalence classes of finite-dimensional central simple \(F\)-algebras, with multiplication induced by tensor product. Two such algebras are Brauer equivalent if suitable matrix algebras over them are isomorphic as \(F\)-algebras.

A cyclic algebra of degree \(4\) over \(F\) is an algebra \((L/F,\sigma,a)\), where \(L/F\) is a cyclic Galois extension of degree \(4\), \(\sigma\) generates \(\operatorname{Gal}(L/F)\), and \(a\in F^\times\). It has underlying vector space \(\bigoplus_{j=0}^{3}Lu^j\), with multiplication specified by \(u^4=a\) and \(u\ell=\sigma(\ell)u\) for \(\ell\in L\). Define
\[
C_4(F)=\left\langle [(L/F,\sigma,a)] : L/F\text{ cyclic Galois of degree }4,\ \sigma\text{ a generator},\ a\in F^\times\right\rangle
\subseteq\operatorname{Br}(F).
\]
Here generation means generation as a subgroup, so its elements are finite products of the indicated classes and their inverses.

The supplied Problem 9 asks whether, for every field \(F\) satisfying
\[
\operatorname{char}(F)\ne 2,\qquad -1\notin(F^\times)^2,
\]
one has
\[
\operatorname{Br}_p[4](F)=C_4(F).
\]
The second hypothesis is the supplied condition \(\sqrt{-1}\notin F\). The notation \(\operatorname{Br}_p[4](F)\) is retained exactly from the record, but neither this notation nor the subscript \(p\) is defined there. Consequently, the intended subgroup on the left remains unresolved, and the displayed question requires clarification before it is a fully specified mathematical statement. In particular, it has not been silently identified with the subgroup of classes annihilated by \(4\).

The source introduces this as a question arising for composite \(n\), with \(n=4\) in the displayed instance. It also supplies the following related question, attributed to M. Mahdavi-Hezavehi: for every field \(F\), does divisibility of \(F^\times\) imply \(\operatorname{Br}(F)=0\)? Divisibility means that for every \(a\in F^\times\) and every integer \(m\geq1\), there exists \(b\in F^\times\) with \(b^m=a\). As background, the source states that this related question has an affirmative answer when \(\operatorname{Br}(F)\) is generated by cyclic algebras, because these algebras split over such a field. It further states that the answer is affirmative in characteristic zero, since divisibility then supplies all roots of unity and the Merkurjev--Suslin theorem applies. These assertions are retained as supplied source qualifications, without independent verification.
\subsection{Short English statement}
For every field \(F\) of characteristic different from \(2\) that does not contain a square root of \(-1\), the source asks whether \(\operatorname{Br}_p[4](F)\) is generated by classes of cyclic algebras of degree \(4\). The target notation is undefined in the supplied record, so its intended meaning requires review. The registry labels the problem open; that label is supplied curation, not an independently verified status assessment.
\subsection{Sources}
\begin{itemize}
\item Local UnsolvedMath registry, permanent ID \(800699\), exact source problem code \texttt{catalogue entry}, titled \emph{Open Problems on Central Simple Algebras --- Problem 9}. The supplied record attributes the problem to Auel, Brussel, Garibaldi, and Vishne, with proposed year \(2010\), and labels it open.
\item Auel, Brussel, Garibaldi, and Vishne, \emph{Open Problems on Central Simple Algebras} (2010), Problem 9. Supplied source URL: \url{https://arxiv.org/abs/1006.3304}. The record reports access on \texttt{2026-10-16} and extraction by \texttt{arxiv\_latex}; these are supplied provenance claims, not claims of access or verification for this draft.
\item The related divisibility question is attributed in the supplied text to M. Mahdavi-Hezavehi. The same text invokes the Merkurjev--Suslin theorem for its characteristic-zero qualification, but supplies no separate bibliographic reference for either attribution.
\end{itemize}
\end{document}
